\begin{abstract}
Mathematical reasoning remains a formidable obstacle for language models, primarily because of its inherent complexity and structured requirements. In this work, we present DeepSeekMath~7B, which extends the pre-training of DeepSeek-Coder-Base-v1.5 7B utilizing 120B math-centric tokens drawn from Common Crawl, along with data from natural language and code domains. Without the assistance of external toolkits or ensembling strategies, DeepSeekMath~7B attains a notable score of 51.7\% on the MATH competition benchmark, approaching the proficiency of Gemini-Ultra and GPT-4. When employing self-consistency across 64 generated samples, DeepSeekMath~7B's performance on MATH increases to 60.9\%. The improvements in mathematical reasoning within DeepSeekMath~7B can be credited to two main innovations: First, a rigorously designed data selection framework that leverages the extensive pool of open web resources; second, the introduction of Group Relative Policy Optimization (GRPO), an adaptation of Proximal Policy Optimization (PPO), which not only strengthens mathematical reasoning capabilities but also streamlines PPO's memory requirements.
\end{abstract}

\section{Introduction}

Large language models (LLMs) have dramatically transformed artificial intelligence approaches to mathematical reasoning, driving notable progress in benchmarks for quantitative reasoning \citep{MATH} and geometric reasoning \citep{trinh2024solving}. These models have also become vital tools in supporting human users with intricate mathematical problem solving \citep{tao}. Yet, state-of-the-art systems like GPT-4 \citep{gpt4} and Gemini-Ultra \citep{gemini} remain inaccessible to the public, and current open-source alternatives lag substantially in their mathematical reasoning performance.

This paper presents DeepSeekMath, a specialized language model tailored for mathematics that demonstrates superior mathematical reasoning compared to existing open-source models and nearly matches GPT-4 on a range of academic benchmarks. Central to this advancement is the construction of the DeepSeekMath Corpus, an extensive and meticulously curated pre-training dataset containing 120B math-specific tokens. The corpus is assembled from the Common Crawl (CC) utilizing a classifier built on fastText \citep{joulin2016fasttext}. Initially, the classifier is trained with positive samples drawn from OpenWebMath \citep{openwebmath} and a broad array of unrelated web pages serving as negative samples. This trained classifier is then deployed to identify additional mathematical content within CC, and these newly identified samples undergo further validation through human annotation. The classifier is iteratively retrained with this expanded and refined dataset, enhancing its selection precision. Our empirical results demonstrate the robustness of the resulting corpus, as evidenced by DeepSeekMath-Base 7B attaining a score of 64.2\% on GSM8K \citep{gsm8k} and 36.2\% on the high-difficulty MATH dataset \citep{MATH}, surpassing the performance of Minerva 540B \citep{minerva}. Furthermore, due to the multilingual nature of the DeepSeekMath Corpus, our model exhibits notable improvements on Chinese mathematical tasks \citep{wei2023cmath,agieval}. We regard our methodology in constructing mathematical data resources as an initial foundation for further research, anticipating considerable opportunities for advancement in this domain.

DeepSeekMath-Base is built upon DeepSeek-Coder-Base-v1.5 7B \citep{deepseek-coder}, since initializing from a code-focused model yields better results than starting with a generic LLM. Notably, incorporating math-specific training also elevates model performance on benchmarks such as MMLU \citep{mmlu} and BBH \citep{bbh}, suggesting that mathematical training not only sharpens math-specific skills but also enhances the model's broader reasoning capacities.

Following pre-training, we perform mathematical instruction tuning on DeepSeekMath-Base by leveraging data involving chain-of-thought reasoning \citep{cot}, program-of-thought processes \citep{pot,pal}, and tool-assisted reasoning techniques \citep{tora}. The resulting model, DeepSeekMath-Instruct 7B, surpasses all other 7B-scale competitors and rivals instruction-tuned models at the 70B scale available in open source.

Additionally, we present Group Relative Policy Optimization (GRPO), a novel reinforcement learning (RL) variant inspired by Proximal Policy Optimization (PPO) \citep{schulman2017proximal}. GRPO eliminates the need for a critic model, opting instead to estimate the baseline using group scores, which leads to considerable savings in computational resources. Utilizing only a fraction of English instruction tuning data, GRPO achieves marked performance gains over the already strong DeepSeekMath-Instruct—covering both in-domain datasets (GSM8K: 82.9\% $\rightarrow$ 88.2\%, MATH: 46.8\% $\rightarrow$ 51.7\%) and out-of-domain problems (for example, CMATH: 84.6\% $\rightarrow$ 88.8\%) throughout the reinforcement learning process. 

We further develop a unified framework for interpreting various methods, including Rejection Sampling Fine-Tuning (RFT) \citep{yuan2023scaling}, Direct Preference Optimization (DPO) \citep{dpo}, PPO, and GRPO. Through this unified lens, we demonstrate that these approaches can all be categorized as direct or streamlined RL algorithms. Moreover, we conduct comprehensive studies—including comparisons between online and offline training, evaluations of outcome versus process-based supervision, and investigations into single-turn versus iterative RL—to thoroughly analyze the core factors of this framework. Finally, we offer insights into why our RL methods enhance the capabilities of instruction-tuned models and outline promising avenues for advancing RL effectiveness within this unified structure.

\subsection{Contributions}
This work advances scalable mathematics pre-training and delves into reinforcement learning through systematic investigation and assessment.

\textbf{Math Pre-Training at Scale}

\begin{itemize}[topsep=0pt]
    \item We demonstrate that the openly available Common Crawl data serves as a rich source for mathematical pre-training. Utilizing a carefully engineered data selection process, we have developed the DeepSeekMath Corpus, a curated dataset comprising 120B tokens derived from web content with mathematical relevance. This corpus surpasses the scale of Minerva’s math dataset \citep{minerva} by nearly sevenfold, and is over nine times larger than OpenWebMath’s recent release \citep{openwebmath}.

    \item Our DeepSeekMath-Base 7B model, trained from scratch, attains performance on par with Minerva 540B \citep{minerva}. This result suggests that the parameter count is not the sole determinant of a model’s prowess in mathematical reasoning. High-quality training data can empower relatively compact models to achieve competitive outcomes.

    \item We present insights gained from our math pre-training experiments. Notably, conducting code pre-training prior to mathematical pre-training enhances models' mathematical problem-solving abilities, with or without the aid of tools. This partially addresses the longstanding debate regarding the impact of code training on reasoning, providing evidence that it benefits mathematical reasoning specifically.

    \item While leveraging arXiv papers for training is a widespread practice in mathematical language modeling, our results reveal no measurable gains on the suite of mathematical benchmarks evaluated in this study.
\end{itemize}

\textbf{Exploration and Analysis of Reinforcement Learning}

\begin{itemize}[topsep=0pt]
    \item We present Group Relative Policy Optimization (GRPO), a novel and resource-efficient reinforcement learning approach. Rather than relying on a critic network, GRPO utilizes group-based scores to estimate the baseline, thereby greatly decreasing the computational demands relative to Proximal Policy Optimization (PPO).
    \item Through experiments, we show that GRPO provides substantial improvements to the instruction-tuned DeepSeekMath-Instruct model, using only instruction-tuning datasets. Additionally, we detect marked improvements in generalization to out-of-domain tasks during reinforcement learning.
    \item We propose a comprehensive framework to contextualize various techniques, such as RFT, DPO, PPO, and GRPO. We further perform thorough analyses, including comparisons between online and offline training, evaluation of outcome-based versus process-based supervision, and studies of single-step versus iterative reinforcement learning, to systematically examine core aspects of the paradigm.
    \item Leveraging this unified framework, we delve into the underlying mechanisms that make reinforcement learning effective, and outline several avenues for advancing the reinforcement learning of large language models (LLMs).
\end{itemize}

\subsection{Summary of Evaluations and Metrics}

\begin{itemize}[topsep=0pt]
    \item \textbf{English and Chinese Mathematical Reasoning}: We carry out detailed evaluations of our models on a wide range of English and Chinese benchmarks, encompassing mathematical problems spanning from elementary to university levels. The English suite includes GSM8K \citep{gsm8k}, MATH \citep{MATH}, SAT \citep{llemma}, OCW Courses \citep{minerva}, and MMLU-STEM \citep{mmlu}; while the Chinese evaluations feature MGSM-zh \citep{mgsm}, CMATH \citep{wei2023cmath}, Gaokao-MathCloze \citep{agieval}, and Gaokao-MathQA \citep{agieval}. Our assessments consider both the capacity of models to produce independent textual solutions (without auxiliary tools) and their problem-solving abilities via Python.

    On English tasks, DeepSeekMath-Base demonstrates strong competitiveness, achieving performance on par with the proprietary Minerva 540B \citep{minerva}, and outperforming all open-source baseline models such as Mistral 7B \citep{mistral} and Llemma-34B \citep{llemma}, regardless of their math pre-training status, often by a considerable margin. DeepSeekMath-Base achieves even greater results on the Chinese benchmarks, a likely consequence of our strategy to include substantial, high-quality multilingual mathematical data, rather than limiting pre-training to English content as done in previous studies \citep{minerva,llemma}. Following further instruction tuning and reinforcement learning, DeepSeekMath-Instruct and DeepSeekMath-RL yield remarkable outcomes, exceeding 50\% accuracy on the challenging, competition-level MATH benchmark—a milestone achieved for the first time among open-source models.
\end{itemize}

\item \textbf{Formal Mathematics}: We assess DeepSeekMath-Base's abilities on the informal-to-formal theorem proving challenge introduced in \citep{dsp_proof}, employing the miniF2F dataset \citep{minif2f} and using Isabelle \citep{isabelle} as the proof assistant. DeepSeekMath-Base exhibits robust few-shot autoformalization capabilities.
\item \textbf{Natural Language Understanding, Reasoning, and Code}: To thoroughly characterize the model's general competencies in language comprehension, reasoning, and programming, we subject DeepSeekMath-Base to the Massive Multitask Language Understanding (MMLU) benchmark \citep{mmlu}, which features 57 diverse multiple-choice subjects, as well as BIG-Bench Hard (BBH) \citep{bbh}, a suite of 23 particularly challenging problems primarily demanding complex, multi-step reasoning. Additionally, we evaluate performance on HumanEval \citep{codex} and MBPP \citep{mbpp}, benchmarks commonly adopted for assessing code generation models. Pre-training on mathematical content is shown to enhance results across both reasoning and natural language understanding benchmarks.

\section{Math Pre-Training}

\subsection{Data Collection and Decontamination}
This section describes the methodology for assembling the DeepSeekMath Corpus from Common Crawl. As illustrated in Figure \ref{fig:data_collect}, we introduce an iterative framework that enables the systematic construction of a large-scale mathematical dataset from Common Crawl, beginning with an initial high-quality “seed” (such as a curated subset of math datasets). The generalizability of this framework extends to other specialized domains, including programming.

Our process begins with OpenWebMath \citep{openwebmath}, a carefully filtered compilation of mathematical texts from the web, as the seed dataset. Leveraging this resource, we train a fastText model \citep{joulin2016fasttext} to identify and extract further web pages similar in mathematical character to those found in OpenWebMath. For model training, 500,000 samples randomly drawn from the seed serve as positive instances, while 500,000 web pages sampled from Common Crawl serve as negatives. The training utilizes the official fastText library\footnote{\url{https://fasttext.cc}}, set to a vector size of 256, a learning rate of 0.1, a maximum n-gram length of 3, a minimum occurrence of 3 words, and three training epochs. In order to scale down the overall Common Crawl dataset, we first apply deduplication and near-deduplication methods based on URLs, which reduces the pool to 40B HTML pages. Next, we use the trained fastText model to retrieve mathematical documents from the deduplicated web pages. To ensure quality, these retrieved pages are scored according to the fastText model’s predictions and only the highest-scoring (i.e., top-quality) pages are retained. The retained dataset size is then determined based on experimental results from pre-training on subsets containing the top 40B, 80B, 120B, and 160B tokens, with the initial iteration preserving the top 40B tokens.

Following the initial round of data collection, a substantial portion of mathematical web pages remains uncollected. This gap is primarily due to the limited diversity in the set of positive examples used to train the fastText model. To address this, we enhance the seed corpus by sourcing additional mathematical web content, thereby improving the fastText model's effectiveness. We begin by partitioning the entire Common Crawl dataset into distinct domains, where each domain comprises web pages with the same base URL. For every domain, we determine the fraction of web pages retrieved during the first round. If more than 10\% of a domain's pages have been collected, we label that domain as math-related (for example, \url{mathoverflow.net}). 

We then conduct a manual review to annotate URLs within these domains that are specifically associated with mathematical material (e.g., \url{mathoverflow.net/questions}). Uncollected web pages that are connected to these URLs are incorporated into the seed corpus. This method allows us to aggregate a broader set of positive samples, facilitating the training of a more robust fastText model that can recover more mathematical content in the next iteration. After conducting four iterations of this process, we accumulate a total of 35.5 million mathematical web pages, corresponding to 120 billion tokens. By the fourth iteration, we observe that almost 98\% of the mathematical pages had already been obtained by the end of the third round, leading us to discontinue further collection.

To safeguard against benchmark leakage, we adopt the protocol outlined in \cite{deepseek-coder} by filtering out web pages that contain material from English mathematical benchmarks such as GSM8K \citep{gsm8k} and MATH \citep{MATH}, as well as Chinese benchmarks like CMATH \citep{wei2023cmath} and AGIEval \citep{agieval}. Specifically, any segment of text that contains a 10-gram sequence identical to any substring from these benchmarks is eliminated from our math training set. For benchmark materials shorter than 10 grams but containing at least 3 grams, we utilize an exact match strategy to remove the contaminated web pages.

\subsection{Validating the Quality of the DeepSeekMath~Corpus}

To evaluate the quality of the DeepSeekMath~Corpus, we perform pre-training experiments comparing it to recently released math-oriented corpora:

\begin{itemize}[topsep=0pt]
    \item \textbf{MathPile} \citep{mathpile}: This multi-source corpus (8.9B tokens) is compiled from textbooks, Wikipedia, ProofWiki, CommonCrawl, StackExchange, and arXiv, with over 85\% of the tokens originating from arXiv.
    \item \textbf{OpenWebMath} \citep{openwebmath}: Derived from CommonCrawl and filtered for mathematical content, encompassing 13.6B tokens.
    \item \textbf{Proof-Pile-2} \citep{llemma}: Comprising OpenWebMath, AlgebraicStack (10.3B tokens of mathematical code), and arXiv articles (28.0B tokens). When using Proof-Pile-2 in experiments, we follow \cite{llemma} and maintain an arXiv:Web:Code ratio of 2:4:1.
\end{itemize}

\subsubsection{Training Setting}
\label{sec:quality-policy}
We adapt a general-purpose pre-trained language model consisting of 1.3B parameters—which maintains architectural compatibility with DeepSeek LLMs \citep{deepseek-llm}, hereafter referred to as DeepSeek-LLM 1.3B—by fine-tuning it on individual mathematical corpora, each over the span of 150B tokens. For all our training procedures, we employ the resource-efficient HAI-LLM framework \citep{haillm}. In accordance with the DeepSeek LLMs training protocols, the AdamW optimizer \citep{adamW} is used, setting $\beta_1=0.9$, $\beta_2=0.95$, and a $\mathrm{weight\_decay}$ of 0.1. The learning rate is managed with a multi-stage schedule: it ramps up to its maximum value over 2,000 warmup steps, then decays to 31.6\% of its peak after 80\% of total steps, and drops further to 10.0\% of the maximum after 90\% of training is complete. The learning rate is capped at 5.3e-4. The batch size consists of 4 million tokens, using a context window of 4,096 tokens.

\subsubsection{Evaluation Results}

\textbf{The DeepSeekMath~Corpus excels in quality, provides extensive multilingual mathematical resources, and surpasses others in scale.}

\begin{itemize}[topsep=0pt]
    \item \textbf{High-quality}:
    We assess downstream capabilities on eight distinct mathematical benchmarks, leveraging few-shot chain-of-thought prompting as in \cite{cot}. Table \ref{tab:corpora_comparison} highlights the clear superiority of the model trained with the DeepSeekMath~Corpus. In addition, Figure \ref{fig:corpora_comparisons} illustrates that this model surpasses the one trained on Proof-Pile-2 even after processing 50B tokens (one complete pass through Proof-Pile-2), signifying that DeepSeekMath Corpus maintains a higher average data quality.
    \item \textbf{Multilingual}:
    DeepSeekMath~Corpus consists of material spanning various languages, with English and Chinese forming the primary language bases. Table \ref{tab:corpora_comparison} shows that models trained on DeepSeekMath~Corpus demonstrate notable gains in mathematical reasoning in both English and Chinese. In comparison, extant mathematical corpora—largely English-dominated—tend not to boost, and sometimes even degrade, mathematical reasoning capabilities in Chinese.
    \item \textbf{Large-scale}:
    DeepSeekMath~Corpus is significantly larger than other available mathematical corpora. As presented in Figure \ref{fig:corpora_comparisons}, DeepSeek-LLM 1.3B achieves a more pronounced improvement in performance and sustains learning for longer periods when trained on DeepSeekMath~Corpus. Conversely, smaller baselines see their training corpora repeated across several epochs, with the corresponding models plateauing in performance relatively quickly.
\end{itemize}

\subsection{Training and Evaluating DeepSeekMath-Base 7B}

Here, we present DeepSeekMath-Base 7B, a foundational model distinguished by its robust reasoning skills, particularly within the mathematical domain. The model initialization leverages DeepSeek-Coder-Base-v1.5 7B \citep{deepseek-coder}, and it undergoes training over 500B tokens. The data utilized in training is sourced as follows: 56\% originates from the DeepSeekMath Corpus, 4\% from AlgebraicStack, 10\% from arXiv, 20\% from code repositories on Github, with the remaining 10\% comprising natural language samples from Common Crawl in English and Chinese. Training procedures mostly adhere to the setup delineated in Section \ref{sec:quality-policy}, with two key exceptions: the peak learning rate is set to 4.2e-4, and the batch size is established at 10M tokens.

To thoroughly evaluate DeepSeekMath-Base 7B, we assess its mathematical proficiency across various dimensions: its capacity to independently generate fully detailed mathematical solutions, to utilize auxiliary tools when tackling problems, and to engage in formal theorem proving. Additionally, the evaluation extends beyond mathematical tasks, encompassing analyses of the base model’s abilities in natural language understanding, logical reasoning, and programming.

\paragraph{Mathematical Problem Solving with Step-by-Step Reasoning}

To measure DeepSeekMath-Base’s aptitude in mathematical problem solving, we implement few-shot chain-of-thought prompting \citep{cot}, examining performance on eight distinct English and Chinese benchmarks. These benchmarks include datasets focused on quantitative reasoning, such as GSM8K \citep{gsm8k}, MATH \citep{MATH}, and CMATH \citep{wei2023cmath}, as well as multiple-choice assessments like MMLU-STEM \citep{mmlu} and Gaokao-MathQA \citep{agieval}. The suite of benchmarks represents a wide array of mathematical domains, spanning elementary topics up to university-level mathematics.

As presented in Table \ref{tab:base_math_cot}, DeepSeekMath-Base 7B achieves the highest results across all eight evaluated benchmarks among open-source base models, including popular choices such as Mistral 7B \citep{mistral} and the newly introduced Llemma 34B \citep{llemma}, the latter of which is specifically trained on mathematical data from Proof-Pile-2 \citep{llemma}. Particularly noteworthy is DeepSeekMath-Base’s performance on the MATH dataset at the competition level, where it exceeds the next-best open-source model by over 10\% in absolute terms, and also outperforms Minerva 540B \citep{minerva}—a substantially larger, closed-source model based on PaLM \citep{palm} and further refined on mathematical corpora.

\paragraph{Mathematical Problem Solving with Tool Use} To assess models’ capabilities in program-aided mathematical reasoning, we employ few-shot program-of-thought prompting \citep{pot,pal} on the GSM8K and MATH datasets. Each prompt requires the model to generate a Python script, utilizing libraries such as \textit{math} and \textit{sympy} to handle complex calculations. The program’s output is then used as the problem’s solution. Table \ref{tab:base_math_pot_proof} reveals that DeepSeekMath-Base 7B outperforms the previous leading model, Llemma 34B.

\paragraph{Formal Mathematics}

Automating formal proofs is crucial for bolstering proof correctness and efficiency, and this task has garnered considerable interest. We evaluate DeepSeekMath-Base 7B on the informal-to-formal proof generation task introduced by \citep{dsp_proof}. In this task, the model is provided an informal statement, its corresponding formal version, and an informal proof, and is asked to generate a formal proof. Evaluations are conducted on the miniF2F \citep{minif2f} dataset, targeting formal mathematical problems at the Olympiad level, with solutions produced in Isabelle using few-shot prompts. Following the procedure in \cite{dsp_proof}, the model generates proof outlines, and we utilize the Sledgehammer automated prover \citep{sledgehammer} to complete any remaining gaps. According to Table \ref{tab:base_math_pot_proof}, DeepSeekMath-Base 7B delivers robust performance in automatic proof formalization.

\paragraph{Natural Language Understanding, Reasoning, and Code}

The model’s abilities in natural language understanding, logical reasoning, and code generation are tested using MMLU \citep{mmlu} for comprehension, BBH \citep{bbh} for reasoning, and HumanEval \citep{codex} together with MBPP \citep{mbpp} for coding. Table \ref{tab:understanding_reasoning_code} indicates that DeepSeekMath-Base 7B substantially improves performance over its predecessor, DeepSeek-Coder-Base-v1.5 \citep{deepseek-coder}, particularly in MMLU and BBH, highlighting the benefits of mathematical pretraining for language comprehension and reasoning. By continuing training with code tokens, DeepSeekMath-Base 7B also preserves the coding performance achieved by DeepSeek-Coder-Base-v1.5 on the programming benchmarks. In sum, DeepSeekMath-Base 7B consistently outshines the general-purpose Mistral 7B \citep{mistral} across all three benchmarks for reasoning and programming.

\section{Supervised Fine-Tuning}

\subsection{SFT Data Curation}

We assembled an instruction-tuning dataset for mathematics that incorporates both English and Chinese questions, sourced from a broad range of mathematical domains and exhibiting diverse levels of difficulty. Each problem is paired with corresponding solutions structured in chain-of-thought (CoT) \citep{cot}, program-of-thought (PoT) \citep{pot,pal}, or tool-integrated reasoning formats \citep{tora}. The resulting dataset comprises 776K training instances.

\begin{itemize}[topsep=0pt]
    \item \textbf{English mathematical datasets}: For English data, GSM8K and MATH datasets were provided with tool-integrated annotated answers. Additionally, we utilized a portion of MathInstruct \citep{MathInstruct} and the training examples from Lila-OOD \citep{lila}, in which problems are solved utilizing either CoT or PoT. This dataset encompasses a wide array of mathematical areas, including but not limited to algebra, probability, number theory, calculus, and geometry.
    \item \textbf{Chinese mathematical datasets}: For Chinese content, we aggregated K-12 mathematics questions from 76 distinct subfields, such as linear equations, and supplied them with both CoT and tool-integrated reasoning solution annotations.
\end{itemize}

\subsection{Training and Evaluating DeepSeekMath-Instruct 7B}

Here, we describe the fine-tuning process for DeepSeekMath-Instruct 7B, which builds upon DeepSeekMath-Base through mathematical instruction-tuning. During training, examples are concatenated in random order until a maximum token limit of 4K is met. The model is trained for a total of 500 steps with a batch size set at 256 and uses a fixed learning rate of 5e-5.

Mathematical capabilities are evaluated in both settings—without tools and with tool assistance—across four quantitative reasoning benchmarks in both English and Chinese. We conduct comparisons with leading contemporary models, outlined below:

\begin{itemize}[topsep=0pt]
    \item \textbf{Closed-source models}: These consist of the GPT series, with GPT-4 \citep{gpt4} and GPT-4 Code Interpreter~\footnote{\url{https://openai.com/blog/chatgpt-plugins\#\#code-interpreter}} being the most advanced; Gemini Ultra and Pro \citep{gemini}; Inflection-2 \citep{inflection-2}; Grok-1~\footnote{\url{https://x.ai/model-card}}; along with recently published models from Chinese technology companies, including Baichuan-3~\footnote{\url{https://www.baichuan-ai.com}} and the latest member of the GLM family, GLM-4~\footnote{\url{https://open.bigmodel.cn/dev/api\#glm-4}} \citep{glm}.
\end{itemize}

These models are designed for broad applicability, and most have been subjected to a series of alignment procedures.

\item \textbf{Open-source models} encompass both general-purpose and math-enhanced architectures: (1) DeepSeek-LLM-Chat 67B \citep{deepseek-llm}, (2) Qwen 72B \citep{qwen}, (3) SeaLLM-v2 7B \citep{seallm}, and (4) ChatGLM3 6B \citep{chatglm3} represent standard general models, while those incorporating mathematical advancements include (5) InternLM2-Math 20B~\footnote{\url{https://github.com/InternLM/InternLM-Math}}, an extension of InternLM2 that receives additional math-specific instruction tuning, (6) Math-Shepherd-Mistral 7B, which integrates PPO optimization \citep{schulman2017proximal} into Mistral 7B \citep{mistral} using a process-supervised reward model, (7) the WizardMath family \citep{wizardmath}, which augments Mistral 7B and Llama-2 70B \citep{llama2} with improved mathematical reasoning by employing evolve-instruct (instruction tuning with AI-generated prompts) and PPO on datasets such as GSM8K and MATH, (8) MetaMath 70B \citep{metamath}, a Llama-2 70B variant further fine-tuned with an expanded GSM8K and MATH dataset, (9) ToRA 34B \cite{tora}, which adapts CodeLlama 34B for tool-supported mathematical reasoning, and (10) MAmmoTH 70B \citep{MathInstruct}, which features Llama-2 70B instruction-tuned on MathInstruct.

Table \ref{tab:sft_rl_math} displays the results for settings in which external tools are not utilized; here, DeepSeekMath-Instruct 7B exhibits robust stepwise reasoning performance. Particularly for the advanced MATH competition dataset, this model outperforms all other open-source options as well as most closed-source systems (including Inflection-2 and Gemini Pro), achieving a minimum advantage of 9\% in absolute terms. This lead holds even over considerably larger models (e.g., Qwen 72B) or those with targeted mathematical reinforcement learning (such as WizardMath-v1.1 7B). DeepSeekMath-Instruct achieves comparable results to proprietary Chinese models like GLM-4 and Baichuan-3 on MATH, though it still lags behind GPT-4 and Gemini Ultra.

When models are permitted to combine natural language reasoning with code-based tool use for solving problems, DeepSeekMath-Instruct 7B nearly attains 60\% accuracy on MATH—outperforming all other available open-source models. On additional benchmarks, this model demonstrates results on par with DeepSeek-LLM-Chat 67B, the preceding state-of-the-art, despite being just a tenth of the size.

\section{Reinforcement Learning}

\subsection{Group Relative Policy Optimization}
Reinforcement learning (RL) is widely recognized as a highly effective strategy for further enhancing the mathematical reasoning abilities of LLMs following Supervised Fine-Tuning (SFT) \citep{wang2023math,wizardmath}. Here, we present Group Relative Policy Optimization (GRPO), a new RL algorithm that achieves both efficiency and effectiveness.

\subsubsection{From PPO to GRPO} Proximal Policy Optimization (PPO) \citep{schulman2017proximal} remains one of the most prevalent actor-critic RL algorithms for fine-tuning LLMs \citep{ouyang2022training}. PPO refines LLMs by seeking to maximize the surrogate objective given by:

\begin{equation}
\footnotesize
    \mathcal{J}_{PPO}(\theta) = \mathbb{E}{[q \sim P(Q), o \sim \pi_{\theta_{old}}(O|q)]} \frac{1}{|o|} \sum_{t=1}^{|o|} \min \left[ \frac{\pi_\theta(o_{t} | q, o_{<t})}{\pi_{\theta_{old}}(o_{t} | q, o_{<t})} A_{t}, \text{clip} \left( \frac{\pi_\theta(o_{t} | q, o_{<t})}{\pi_{\theta_{old}}(o_{t} | q, o_{<t})}, 1 - \varepsilon, 1 + \varepsilon \

Here, $\pi_{\theta}$ and $\pi_{\theta_{old}}$ denote the present and preceding policy models, while $q$ and $o$ refer to questions and outputs sampled respectively from the dataset and the old policy $\pi_{\theta_{old}}$. The parameter $\varepsilon$ serves as a clipping hyper-parameter in PPO to stabilize training. The advantage term $A_t$ is estimated using Generalized Advantage Estimation (GAE) \citep{gae}, which relies on the sequence of rewards $\{r_{\ge t}\}$ and an associated value function $V_{\psi}$. Consequently, PPO requires the concurrent training of a value function along with the policy. In order to limit excessive optimization against the reward model, the conventional solution appends a KL divergence penalty, measured per token between the policy and a reference model, to the reward at each token step \citep{ouyang2022training}:

\begin{equation}
    r_{t} = r_\varphi(q, o_{\le t}) - \beta \log\frac{\pi_{\theta}(o_{t}|q, o_{<t})}{\pi_{ref}(o_{t}|q, o_{<t})},
\label{eq:PPO-reward}
\end{equation}

Here, $r_\varphi$ represents the reward model, and $\pi_{ref}$ is the reference (typically initial SFT) model. The KL penalty's scaling factor is given by $\beta$.

Utilizing a value function in PPO often entails maintaining another neural network model on a scale similar to the policy itself, introducing notable demands in terms of memory and computation. Moreover, during reinforcement learning, the value function operates as a baseline to reduce variance in the computed advantages. For large language models (LLMs), however, it is common practice for only the final token to receive a reward from the reward model, creating challenges in training a value function that provides meaningful predictions for every token. To overcome this, we introduce Group Relative Policy Optimization (GRPO), depicted in Figure \ref{fig:grpo}, which dispenses with the need for explicit value function approximation as is used in PPO. Instead, GRPO establishes a baseline by averaging the rewards from multiple samples—outputs produced by the old policy in response to a given question. Concretely, for each question $q$, a set $\{o_1, o_2, \cdots, o_G\}$ of outputs are drawn from $\pi_{\theta_{old}}$, and the policy update maximizes:

\begin{equation}
\footnotesize
\begin{split}
    \mathcal{J}_{GRPO}(\theta) &= \mathbb{E}{[q \sim P(Q), \{o_i\}_{i=1}^G \sim \pi_{\theta_{old}}(O|q)]}  \\
    & \frac{1}{G}\sum_{i=1}^G\frac{1}{|o_i|} \sum_{t=1}^{|o_i|} \left\{ \min \left[ \frac{\pi_\theta(o_{i,t} | q, o_{i,<t})}{\pi_{\theta_{old}}(o_{i,t} | q, o_{i,<t})} \hat{A}_{i,t}, \text{clip} \left( \frac{\pi_\theta(o_{i,t} | q, o_{i,<t})}{\pi_{\theta_{old}}(o_{i,t} | q, o_{i,<t})}, 1 - \varepsilon, 1 + \varepsilon \right)  \hat{A}_{i,t} \right] - \beta \mathbb{D}_{KL}\left[\pi_{\theta} || \pi_{ref}\right]\right\} ,
\end{split}
\label{eq:GRPO-obj}
\end{equation}

Here, $\varepsilon$ and $\beta$ continue to serve as hyper-parameters. The advantage term $\hat{A}_{i,t}$ is specifically computed from the relative rewards among outputs within the same group, with its formulation detailed in the next subsections. This approach to computing advantage through relative comparisons among group outputs is well matched to the typical construction of reward models, which are frequently trained on comparative datasets containing responses to identical questions. Furthermore, GRPO eschews embedding the KL penalty into the reward; instead, it regularizes directly through an explicit KL divergence term between the current policy and reference policy in the loss, thereby maintaining the clarity of the calculation for $\hat{A}_{i,t}$. Unlike the per-token KL penalty in (\ref{eq:PPO-reward}), we approximate the KL divergence via the unbiased estimator

\begin{algorithm}[t]
  \small
  \caption{Iterative Group Relative Policy Optimization}
  \textbf{Input} Initial policy model $\pi_{\theta_{\text{init}}}$; reward models $r_\varphi$; task prompts $\mathcal{D}$; 
  hyperparameters $\varepsilon$, $\beta$, $\mu$
  \begin{algorithmic}[1]
    \State Assign policy model: $\pi_\theta \leftarrow \pi_{\theta_{\text{init}}}$
    \For{iteration = 1, \dots, I}
       \State Set reference model: $\pi_{ref} \leftarrow \pi_{\theta}$
      \For{step = 1, \dots, M}
      \State Randomly draw a batch $\mathcal{D}_b$ from $\mathcal{D}$
      \State Synchronize old policy model: $\pi_{\theta_{old}} \leftarrow \pi_{\theta}$ 
      \State For each prompt $q \in \mathcal{D}_b$, generate $G$ outputs $\{o_i\}_{i=1}^G$ by sampling from $\pi_{\theta_{old}} (\cdot \mid q) $
      \State Evaluate outputs with reward model $r_{\varphi}$ to obtain $\{r_i\}_{i=1}^{G}$
      \State Estimate group relative advantage $\hat{A}_{i,t}$ for the $t$-th token in output $o_i$
      \For{GRPO iteration = 1, \dots, $\mu$}
        \State Optimize policy model $\pi_{\theta}$ by maximizing the GRPO objective (Equation \ref{eq:GC-GRPO})
      \EndFor
    \EndFor 
    \State Update $r_\varphi$ using ongoing training with a replay buffer. 
    \EndFor 
  \end{algorithmic}
  \textbf{Output} $\pi_\theta$
  \label{alg:iter-grpo}
\end{algorithm}

\subsubsection{Outcome Supervision RL with GRPO} More specifically, for each question $q$, the old policy $\pi_{\theta_{old}}$ is used to sample a collection of outputs $\{o_1, o_2, \cdots, o_G\}$. Each output is evaluated by a reward model, resulting in $G$ scores $\mathbf{r}=\{r_1, r_2, \cdots, r_G\}$. These reward values are standardized by subtracting their mean and dividing by the standard deviation within the group. With outcome supervision, the normalized reward at the final token of every output $o_i$ is assigned, and all tokens in $o_i$ receive this value as their advantage, i.e., $\hat{A}_{i, t} = \widetilde{r}_i = \frac{r_i- {\rm mean}(\mathbf{r})}{{\rm std}(\mathbf{r})}$. The policy is subsequently updated by maximizing the objective given in equation (\ref{eq:GRPO-obj}).

\subsubsection{Process Supervision RL with GRPO} Since outcome supervision delivers reward signals only at the conclusion of each output—which may inadequately supervise the model for intricate mathematical reasoning—we additionally consider process supervision as described in \cite{wang2023math}. Process supervision assigns a reward to every reasoning step within each generated output. Specifically, given question $q$ and the sampled responses $\{o_1, o_2, \cdots, o_G\}$, a step-level reward model provides per-step evaluations, producing a set of scores: $\mathbf{R} = \{ \{r_1^{index(1)},\cdots,r_1^{index(K_1)}\}, \cdots,  \{r_G^{index(1)},\cdots,r_G^{index(K_G)}\} \}$, where $index(j)$ marks the terminal token index for step $j$ and $K_i$ represents the number of reasoning steps in $o_i$. Each reward is normalized against the mean and standard deviation of $\mathbf{R}$: $\widetilde{r}_i^{index(j)} = \frac{r_i^{index(j)} - {\rm mean(\mathbf{R})}}{{\rm std(\mathbf{R})}}$. The advantage for a given token is then the cumulative sum of normalized rewards assigned to steps whose end indices follow or coincide with that token, i.e., $\hat{A}_{i, t} = \sum_{index(j) \ge t} \widetilde{r}_i^{index(j)}$. This advantage is used to maximize the policy objective as indicated

\subsubsection{Iterative RL with GRPO} During reinforcement learning, the previously trained reward model may become inadequate for effectively supervising the continually improving policy model. To address this, we investigate iterative RL utilizing GRPO. As illustrated in Algorithm \ref{alg:iter-grpo}, our iterative GRPO framework involves regenerating training data for the reward model by leveraging samples produced by the latest policy model. The reward model is retrained in each iteration, employing a replay buffer that includes 10\% of past data. Subsequently, we designate the updated policy model as the new reference and proceed to further optimize the policy model under the supervision of the retrained reward model.

\subsection{Training and Evaluating DeepSeekMath-RL}

Our RL procedure is implemented on the DeepSeekMath-Instruct 7B architecture. The RL dataset comprises approximately 144K chain-of-thought questions focused on GSM8K and MATH, sourced from the SFT training data. We omit SFT questions unrelated to these datasets to more precisely assess the effects of RL on benchmarks for which the RL phase provides no direct supervision. To build the reward model training set, we follow the methodology in \citep{wang2023math}. The initial reward model is trained from DeepSeekMath-Base 7B, adopting a learning rate of 2e-5. For GRPO, the learning rate for updating the policy model is 1e-6, with a KL regularization coefficient of 0.04. Each training instance yields $64$ sampled responses. We use a maximum sequence length of 1024 tokens and set the batch size to 1024. The policy model is updated once per exploration cycle. The evaluation protocol for DeepSeekMath-RL 7B is consistent with that for DeepSeekMath-Instruct 7B: GSM8K and MATH, when formulated as chain-of-thought tasks, are treated as in-domain, while all other benchmarks are categorized as out-of-domain.

Table \ref{tab:sft_rl_math} presents a comparison of open- and closed-source models employing both chain-of-thought and tool-integrated reasoning on English and Chinese benchmark datasets. Our observations indicate the following: 1) DeepSeekMath-RL 7B achieves accuracy rates of 88.2\% on GSM8K and 51.7\% on MATH through the use of chain-of-thought reasoning. These results not only surpass those of all open-source models within the 7B to 70B scale but also exceed the performance of most closed-source models. 2) Importantly, DeepSeekMath-RL 7B's training was limited to chain-of-thought-style instruction tuning on GSM8K and MATH datasets, initializing from DeepSeekMath-Instruct 7B. Even with this narrow range of training data, DeepSeekMath-RL 7B consistently outperforms DeepSeekMath-Instruct 7B across all metrics, underlining the advantages conferred by reinforcement learning.

\section{Discussion}
In the following section, we elaborate on our empirical insights derived from both pre-training and RL phases.

\subsection{Lessons Learnt in Pre-Training}

We begin by detailing our insights from the pre-training phase. Unless explicitly mentioned, all descriptions refer to the experimental protocols described in Section \ref{sec:quality-policy}. Additionally, it should be emphasized that references to the DeepSeekMath Corpus in this discussion pertain to an 89B-token dataset assembled during the second cycle of data curation.

\subsubsection{Code Training Benefits Mathematical Reasoning}

It has been widely speculated—though with limited empirical verification—that exposure to code during training bolsters reasoning abilities. We contribute evidence within the context of mathematical reasoning, indicating that incorporating code data in training facilitates improvement in mathematical reasoning tasks, regardless of tool assistance.

To evaluate the effect of code training on mathematical reasoning, we explored both two-stage training and one-stage training regimens:

\textbf{Two-Stage Training}

\begin{itemize}[topsep=0pt]
    \item \textbf{Code Training for 400B Tokens $\rightarrow$ Math Training for 150B Tokens}: DeepSeek-LLM 1.3B is first pre-trained on 400B code tokens and then further trained with 150B tokens consisting of mathematical data.
    \item \textbf{General Training for 400B Tokens $\rightarrow$ Math Training for 150B Tokens}: For comparison, we substitute the code tokens in the initial stage with general tokens, sampled from an extensive, general-purpose corpus assembled by DeepSeek-AI. This baseline is designed to clarify whether pre-training with code offers unique benefits for mathematical reasoning compared to general text.
\end{itemize}

\textbf{One-Stage Training}

\begin{itemize}[topsep=0pt]
    \item \textbf{Math Training for 150B Tokens}: In this approach, DeepSeek-LLM 1.3B is exposed to 150B tokens containing only mathematical content.
    \item \textbf{Training on a mixture of 400B Code Tokens and 150B Math Tokens}: Because sequential math training after code training leads to diminished coding capabilities, we additionally explore joint training on a combined corpus (400B code tokens + 150B math tokens) in a single stage, seeking both improved mathematical reasoning and preservation of coding ability.
\end{itemize}

\paragraph{Results}
The findings summarized in Table \ref{tab:code-math} and Table \ref{tab:code-math-general-eval} report how downstream tasks are affected under various training regimes.

Pre-training with code contributes positively to program-aided mathematical reasoning, whether the procedure follows a two-stage or a one-stage training design. Evidence in Table \ref{tab:code-math} indicates that code-centric pre-training alone confers a marked advantage in handling GSM8K and MATH datasets using Python. The inclusion of additional math-specific training further boosts performance. Notably, in the one-stage joint training regime, the concurrent use of both code and math tokens not only counters the catastrophic forgetting observed in two-stage setups but also yields concurrent improvements for both code-based evaluation (Table \ref{tab:code-math-general-eval}) and mathematical reasoning tasks supported by programmatic tools (Table \ref{tab:code-math}).

Code token training enhances mathematical reasoning capabilities even when external tools are absent. Within a two-stage design, the code-focused pre-training phase alone produces noticeable, though less pronounced, gains. It also increases the efficacy of subsequent math-centric training, culminating in optimal task performance. Conversely, blending code and math tokens in a single-stage setup is less effective for tool-free mathematical reasoning, possibly because DeepSeek-LLM 1.3B, given its restricted parameter count, cannot robustly encode both domains simultaneously.

\subsubsection{ArXiv Papers Appear Ineffectual in Enhancing Mathematical Reasoning}

While arXiv papers are often utilized as a significant component in mathematical pre-training corpora \citep{minerva,gpt-f,llemma,mathpile}, systematic evaluations of their efficacy for fostering mathematical reasoning remain limited. Contrary to common assumptions, our findings indicate that arXiv content has minimal positive influence on models' mathematical reasoning ability. We evaluated models across different parameter scales—specifically, DeepSeek-LLM 1.3B and DeepSeek-Coder-Base-v1.5 7B \citep{deepseek-coder}—by training on arXiv-derived corpora subjected to distinct preprocessing approaches:

\begin{itemize}[topsep=0pt]
    \item \textbf{MathPile} \citep{mathpile}: This is a dataset comprising 8.9B tokens, cleaned and filtered using various heuristic techniques, of which more than 85\% are sourced from scientific arXiv papers.
    \item \textbf{ArXiv-RedPajama} \citep{redpajama}: Consisting of all arXiv LaTeX files with elements such as preambles, comments, macros, and bibliographies excised, accumulating to 28.0B tokens.
\end{itemize}

We carried out separate pre-training for DeepSeek-LLM 1.3B (over 150B tokens) and DeepSeek-Coder-Base-v1.5 7B (over 40B tokens) on each of the arXiv-based corpora. Our results consistently show that training exclusively on arXiv datasets does not bolster, and occasionally hinders, performance on a spectrum of mathematical reasoning benchmarks. These evaluation sets span quantitative tasks such as GSM8K and MATH (see Table \ref{tab:arxiv-cot}), multiple-choice assessments like MMLU-STEM (refer to Table \ref{tab:arxiv-cot}), and formal mathematics challenges exemplified by miniF2F (see Table \ref{tab:arxiv_atp}).

Nevertheless, these results are circumscribed by certain caveats. We did not yet investigate:

\begin{itemize}[topsep=0pt]
    \item The influence of arXiv pre-training data on math-related objectives not examined here, such as translating formal proofs into informal discourse (informalization of theorems);
    \item The effects arising when arXiv data is intermixed with other training sources;
    \item Whether potential gains from arXiv corpora become apparent when scaling to much larger models.
\end{itemize}

Accordingly, more comprehensive investigations are warranted, which we propose as directions for subsequent research.

\subsection{Insights into Reinforcement Learning}
\subsubsection{Toward a Unified Training Framework} This subsection articulates a generalized framework to examine a range of training strategies, including SFT, RFT, DPO, PPO, GRPO, and systematically investigates the parameters influencing these approaches. In a broad sense, the gradient of the training objective with respect to the model parameter $\theta$ is expressible as follows:

\begin{equation}
    \nabla_{\theta}\mathcal{J}_{\textcolor{red}{\mathcal{A}}}(\theta) = \mathbb{E}[\underbrace{(q,o) \sim \textcolor{red}{\mathcal{D}}}_{Data \ Source}]\left( \frac{1}{|o|} \sum_{t=1}^{|o|}  \underbrace{GC_{{\mathcal{A}}}(q, o, t, \textcolor{red}{\pi_{{rf}}})}_{Gradient \ Coefficient}  \nabla_{\theta}\log \pi_{\theta}(o_t | q, o_{<t})\right).
\label{eq:objective}
\end{equation}

This framework is characterized by three principal elements: 1) \textit{Data Source $\mathcal{D}$}, which specifies the origin of the examples for training; 2) \textit{Reward Function $\pi_{{rf}}$}, responsible for generating the reward signal guiding the optimization; and 3) \textit{Algorithm $\mathcal{A}$}, which combines training data and reward information to produce the gradient coefficient $GC$ that modulates the strength and direction of the gradient updates (reinforcement or penalization) for the given data instance. Employing this unified viewpoint, we compare several prototypical approaches:

\begin{itemize}
    \item \textbf{Supervised Fine-tuning (SFT)}: In SFT, the pretrained model is adapted via supervised learning on data curated by human annotators.
    \item \textbf{Rejection Sampling Fine-tuning (RFT)}: Here, the SFT-tuned model is subjected to additional fine-tuning using responses sampled from itself on SFT prompts; these responses are then filtered, retaining only those with correct answers.
    \item \textbf{Direct Preference Optimization (DPO)}: DPO advances the SFT model by further fine-tuning it with pairs of model-generated responses, optimizing a pairwise DPO objective.
    \item \textbf{Online Rejection Sampling Fine-tuning (Online RFT)}: Unlike standard RFT, Online RFT continuously samples augmented responses from the currently updated policy model and employs these for ongoing fine-tuning, starting from the SFT-initialized policy.
    \item \textbf{PPO/GRPO}: With PPO/GRPO, training begins with the SFT model, and model responses sampled in real time are reinforced, refining the policy iteratively.
\end{itemize}

Table \ref{tab:data_policy} presents a summary of the constituent elements of these methods. For an in-depth explanation of the derivation steps, consult Appendix \ref{app:analysis-rl}.

\paragraph{Observation about Data Source} We classify the data source into two primary types: online sampling and offline sampling. Online sampling refers to data generated during the exploration of the actively updated policy model throughout training, whereas offline sampling involves data derived from outputs of the initial SFT model. Among the examined methods, RFT and DPO utilize an offline approach, while Online RFT and GRPO are aligned with the online approach.

From Figure \ref{fig:rl-analysis}, it is evident that Online RFT consistently surpasses RFT across both benchmarks. In the early phases of training, their performances are closely matched; however, as training progresses, Online RFT distinctly outpaces RFT, highlighting the benefits of online training strategies. This trend is plausible, given that early on, the actor closely mirrors the SFT model and thus the differences in sampled data are minimal. At later stages, however, actor-generated samples begin to diverge more substantially from the SFT, and the use of live data sampling becomes markedly advantageous.

\paragraph{Observation about Gradient Coefficient} The algorithms convert the reward signal into a gradient coefficient, which guides parameter updates. In our experiments, the reward function takes two forms: `Rule' and `Model.' The `Rule' approach assesses response quality directly via answer correctness, whereas the `Model' approach employs a trained reward model to evaluate responses, with the reward model being supervised using rule-based labels. Equations \ref{eq:GC-RFT} and \ref{eq:GC-GRPO} illuminate an essential distinction between GRPO and Online RFT: GRPO dynamically modulates its gradient coefficient in accordance with the specific reward model scores, thereby offering granular reinforcement or penalization based on reward magnitude. Conversely, Online RFT lacks such adaptive penalization; it reinforces all correctly answered responses uniformly and does not explicitly penalize incorrect ones.

As depicted in Figure \ref{fig:rl-analysis}, GRPO outperforms online RFT, which emphasizes the effectiveness of adjusting the coefficients for positive and negative gradients. Moreover, the GRPO+PS approach achieves better results than GRPO+OS, suggesting the advantages of utilizing fine-grained, step-aware gradient coefficients. We also investigate the impact of iterative reinforcement learning, conducting two consecutive iterations in our study. Figure \ref{fig:iter-rl} illustrates that iterative RL delivers marked gains, particularly noticeable after the initial iteration.

\subsubsection{Why RL Works?} Our work employs reinforcement learning on a selected subset of instruction tuning data, resulting in notable performance gains over the base instruction-tuned model. To shed more light on the mechanisms underlying the effectiveness of reinforcement learning, we assess the Pass@K and Maj@K metrics for both the Instruct and RL-enhanced models across two benchmark tasks. Figure \ref{fig:maj-pass} demonstrates that RL increases Maj@K accuracy, while Pass@K remains unaffected. These observations suggest that reinforcement learning principally bolsters the robustness of the output distribution—\textbf{the observed improvements are mainly due to elevating correct answers into the TopK outputs rather than an inherent advancement of core capabilities.} Along these lines, \citep{wang2023making} reported a \textbf{misalignment problem} in the context of reasoning tasks tackled by SFT models, indicating that performance can be enhanced via preference alignment approaches \citep{yuan2023rrhf,song2023preference,wang2023making}.

\subsubsection{How to Achieve More Effective RL?}
\label{sec:effec_rl}
Our experiments establish that RL is highly effective on mathematical reasoning challenges. We further offer a unified framework that contextualizes a variety of well-known training methods, portraying them as either direct implementations of, or approximations to, RL strategies. Equation \ref{eq:objective} delineates this paradigm into three fundamental components: Data Source, Algorithm, and Reward Function. We propose several avenues for future research that address each of these aspects.

\paragraph{Data Source} The data source forms the foundational input for all training schemes. For RL, we define this as the set of unlabeled questions and corresponding outputs generated by the policy model. In the current study, we limit ourselves to questions originating from the instruction tuning phase and apply basic nucleus sampling for output generation. This might explain why improvements are primarily seen in Maj@K performance. For future work, we intend to extend the RL pipeline to include out-of-distribution prompts and integrate \textbf{advanced sampling (decoding) strategies}, such as tree-search-based methods \citep{yao2023tree}. In addition, the adoption of \textbf{efficient inference techniques} \citep{xia-etal-2023-speculative,leviathan2023fast,kwon2023efficient,xia2024unlocking}, which directly influence the effectiveness of policy model exploration, are anticipated to play a vital role.

\paragraph{Algorithms} Algorithms use both data and the reward signal to compute the gradient, which in turn updates the model parameters. According to Equation \ref{eq:objective}, current approaches predominantly \textbf{TRUST} the information provided by the reward function to adjust the conditional probability for generating particular tokens. Nevertheless, the reliability of the reward signal cannot always be guaranteed, especially for highly intricate tasks. As an illustration, even the PRM800K datasets \citep{lightman2023let}, which are curated by expert annotators, still display an error rate of roughly 20\% in annotations\footnote{\url{https://github.com/openai/prm800k/issues/12\#issuecomment-1728491852}}. Therefore, investigating reinforcement learning algorithms that are inherently resilient to noisy or erroneous reward signals becomes essential. We anticipate that \textbf{WEAK-TO-STRONG} alignment techniques \citep{burns2023weak} have the potential to fundamentally transform current learning algorithms.

\paragraph{Reward Function} The reward function delivers the core training signal within the RL paradigm, where it typically takes the form of a neural reward model. We identify three pivotal research avenues for the development of reward models: 1) \textbf{Improving the generalization capability of reward models.} It is critical that the reward model robustly generalizes to out-of-distribution queries and sophisticated decoding strategies; otherwise, RL might only maintain the status quo of LLM distributions, failing to enhance their deeper abilities; 2) \textbf{Capturing and communicating the reward model’s uncertainty.} The uncertainty metric may serve as a crucial interface between a limited reward model and the weak-to-strong learning framework; 3) \textbf{Developing effective high-quality process reward models} that yield fine-grained signals to guide the step-by-step reasoning processes in complex tasks \citep{lightman2023let,wang2023math}.

\section{Conclusion, Limitation, and Future Work} In this work, we introduce DeepSeekMath, an open-source model that establishes new state-of-the-art results on the MATH benchmark among openly available systems, coming close to the level of proprietary models. The model is initialized from DeepSeek-Coder-v1.5 7B and receives continuous training on 500B tokens, with a notable proportion (120B tokens) comprising mathematical content extracted from Common Crawl. Through comprehensive ablation analyses, we demonstrate that web pages are a highly valuable resource for mathematical data, while arXiv’s contribution is more limited than anticipated. We propose Group Relative Policy Optimization (GRPO), a novel adaptation of Proximal Policy Optimization (PPO), which significantly enhances mathematical reasoning with reduced memory requirements. Empirical findings indicate that GRPO is effective, even as DeepSeekMath-Instruct 7B attains strong benchmark outcomes. Moreover, we unify several related approaches within a single framework and outline various promising pathways for advancing reinforcement learning effectiveness.

Despite these advancements, DeepSeekMath~exhibits comparatively weaker performance on tasks involving geometry and formal proof than closed-source models. For example, preliminary testing revealed the model’s inability to solve certain triangle and ellipse problems, hinting at potential pre-training and fine-tuning data biases. Furthermore, due to model size limitations, DeepSeekMath~lags behind GPT-4 in few-shot generalization; while GPT-4 benefits from additional examples, DeepSeekMath~performs similarly under both zero-shot and few-shot evaluations. As part of ongoing work, we intend to refine our data selection strategies to curate an even higher-quality pre-training corpus. Additionally, we aim to further investigate the approaches described in Section \ref{sec:effec_rl} for enhancing reinforcement learning in large language models.

\end{document}